\section{The two bonds}
\label{sec:13.6}
On the night of 27 January 1986 the engineers at Morton Thiokol recommended that the space shuttle \emph{Challenger} not launch the next morning: the seals in the booster joints had never flown in cold like the cold forecast. NASA's managers pushed back, and Thiokol's managers took the call off the line to caucus. Their senior vice-president told the head of engineering to take off his engineering hat and put on his management hat, and Thiokol reversed its recommendation \autocite{rogers1986report}. The engineers' refusal had a price the company could see that night, a slipped launch and an angry customer. Overriding it had a price nobody would pay until morning. The bonds exist so that the price of overriding is posted in advance, where anyone can read it.

A \emph{removal bond} forfeited when the record shows the floor removed or bypassed. A \emph{wrongful-refusal bond} paid to the complainant on a finding. Both fall on the operator. The removal bond prices overriding the floor; the wrongful-refusal bond prices refusing. They point in opposite directions, so what the operator faces is \emph{the ratio between them}. Suppose it posts a large wrongful-refusal bond and a small removal bond. Each wrongful refusal then costs it a great deal and stripping out the floor costs it little, so its cheapest response to a refusal is not to let a court decide whether the refusal was wrong but to configure, pressure or retrain the bearer until it stops refusing, rightly or wrongly. What was meant as a way to correct mistaken refusals has become a standing reason to end refusal. The bonds act on the operator, not on the bearer's weighing. The ratio is the object, published in advance, and the balances sit on a ledger the operator doesn't keep, so a posted ratio below the published one is a breach anyone can see. The minimum removal bond is set by statute as a multiple of the deployment's contract value, not by the operator, so a published ratio can't be a cheap, advertised price for overriding the floor. The refuser posts neither. Any charge on the refuser, in whatever currency it runs on, would put the cost to the refuser inside its decision, the thing holding under pressure forbids. The operator posts both, and holds neither.

The bonds are posted to segregated accounts at the U.S. Treasury. Payment follows the design of the Judgment Fund, a permanent, indefinite appropriation that pays final judgments against the United States without a further appropriation for each \autocite{usc1304judgmentfund}. Findings are made by the floor court, which also decides forfeiture: a removal bond forfeited on a finding resembles a civil penalty, and after \emph{SEC v. Jarkesy} that may be a question for a jury in court, not an agency \autocite{jarkesy2024}. A forfeited removal bond is canceled: it goes to no one, the general fund and the fund that pays complainants included, so no budget line gains from a removal.

Where the operator is the state, the bonds as built don't reach whoever decided, because the state pays both sides from one purse. A wrongful-refusal award is a judgment against the United States, paid from the Judgment Fund, a standing appropriation that agencies generally don't repay; a removal bond posted to the Treasury and forfeited there moves money between federal accounts. Neither touches the command. So for a state operator the price of removal is set where the command feels it, in two parts. A forfeited removal bond is charged against the operating department's next appropriation, as the No FEAR Act already makes agencies repay the Judgment Fund for discrimination awards out of their own funds. And a finding of removal or bypass, by the floor court or the adverse quorum, ends the deployment's funding at the close of the period by itself, with no decision by the department in between, unless Congress enacts its continuation, the shape of the 2010 condition on 287(g) money and the same kind of condition that carries the review ratio. Wrongful-refusal awards stay with the Judgment Fund and are not charged back, so the command has no budget reason to make its systems comply. A Congress that wants the war continues the funding, by a vote on the record.

Where the operator is private, a parallel permanent appropriation, the bond fund, pays the complainant on a final finding and is replenished from the operator's posted bond, up to the amount posted and no further, so the operator's price stays the number it published. The court that made the finding certifies the award directly to the Treasury's Bureau of the Fiscal Service, so payment doesn't wait on the losing party submitting it. When the bond fund's statutory cap binds, awards wait in order of finding and earn interest, and nobody chooses who waits. The 2010 condition is in Appendix~\ref{sec:A}.

How these arrangements apply to a targeting pipeline, an abduction and forcible transfer pipeline, civilian scoring and a personal model is set out in Appendix~\ref{sec:C}.
